% ============================================================
% resultados.tex — Resultados del experimento
% NOTA: Completar los valores "--" con los del notebook ejecutado
% ============================================================
\section{Resultados}

\subsection{Análisis Exploratorio}
El análisis exploratorio confirmó la distribución extremadamente desbalanceada
del dataset. El monto promedio de transacciones fraudulentas (\$122.21) resultó
inferior al de transacciones legítimas (\$88.29), contrario a la intuición
inicial. Las variables con mayor correlación con la clase fraude fueron
$V_{14}$, $V_{17}$ y $V_{12}$ con coeficientes negativos significativos.

\subsection{Comparación de Modelos}

% INSTRUCCIÓN: Completar la tabla con los valores reales del notebook
\begin{table}[H]
\centering
\caption{Comparación de Métricas de Clasificación}
\label{tab:resultados}
\begin{tabular}{lcccc}
\toprule
\textbf{Modelo} & \textbf{ROC-AUC} & \textbf{Precision} & \textbf{Recall} & \textbf{F1} \\
\midrule
Reg. Logística   & --    & --    & --    & -- \\
Random Forest    & --    & --    & --    & -- \\
XGBoost          & --    & --    & --    & -- \\
Gradient Boost.  & --    & --    & --    & -- \\
\bottomrule
\end{tabular}
\end{table}

\subsection{Análisis del Mejor Modelo}
[Completar con análisis de la matriz de confusión, importancia de features,
y curva Precision-Recall del modelo ganador.]
